% Source: https://artofproblemsolving.com/community/c3h1370831_help_on_mathcounts_nats_2016_target_problems
% Stable identity: aops-topic-1370831-q05
Take the vertical cross-section through the centre $O$ of the sphere. The rim of the crater is a chord $PQ$ of length $18$ whose midpoint $M$ is $2.5$ inches above the lowest point of the sphere, so with radius $r$ we have $OM=r-2.5$ and $MP=9$.

\begin{asy}
usepackage("amsmath");
size(7cm);
defaultpen(fontsize(10pt));
real R=17.45;
pair O=(0,0), M=(0,-(R-2.5)), P=(9,-(R-2.5)), Q=(-9,-(R-2.5)), Bt=(0,-R);
draw((-22,M.y)--(22,M.y),linewidth(0.9));
draw(circle(O,R),blue);
draw(O--P,red); draw(O--M,red+dashed); draw(M--Bt,gray+dashed);
dot(O);dot(M);dot(P);dot(Q);
label("$O$",O,N); label("$M$",M,NW); label("$P$",P,SE); label("$Q$",Q,SW);
label("$9$",(M+P)/2+(0,1.4),fontsize(9pt)); label("$r$",(O+P)/2,NE,red);
label("$r-2.5$",(O+M)/2,W,red+fontsize(9pt)); label("$2.5$",Bt,S,gray+fontsize(8pt));
label("plaster surface",(-21,M.y),S,fontsize(8pt));
label("\mbox{}",(min(currentpicture,user=true)+max(currentpicture,user=true))/2,black); // keeps the text after the figure black when LaTeX typesets the labels inline
\end{asy}

By the Pythagorean theorem in triangle $OMP$,
\[
r^2=9^2+(r-2.5)^2=81+r^2-5r+6.25\iff5r=87.25\iff r=\boxed{17.45}\text{ inches}.
\]
